\documentclass[12 pt, letterpaper]{hmcpset}
\usepackage{hw-preamble}
\setlist[enumerate, 1]{label = (\roman*)}

\name{Jayden Thadani}
\class{MATH 145 --- Topics in Geometry and Topology}
\assignment{Homework 1 redo}
\duedate{29th January 2025}

\begin{document}

\begin{problem}[1.]
	The following picture is my attempt to draw a $1$-cycle $\alpha$, shown in blue.
	\begin{enumerate}
		\item In fact, I have made mistakes, Specifically, there are exactly three arrows pointing in the wrong direction. Can you identify with certainty which three arrows have been reversed?
		\item Having corrected the errors, compute $w(\alpha, \mathbf{x})$ on each region of $\mathbb{R}^{2} \setminus \lvert \alpha \rvert$.
	\end{enumerate}
\end{problem}

\begin{solution}
	\begin{figure}[H]
		\centering
		\includegraphics[width = 0.72 \textwidth]{figures/P1 cycle.pdf}
		\caption{The cycle $\alpha$ with corrected arrows in red and winding number written inside each region in blue}
	\end{figure}
\end{solution}

\pagebreak

\begin{problem}[2.]
	\begin{enumerate}
		\item Let $\gamma = \sum_{i = 1}^{n} [\mathbf{a}_{i}, \mathbf{b}_{i}]$ be a $1$-cycle, and let $\mathbf{x} \in \mathbb{R}^{2} \setminus \lvert \gamma \rvert$. Show that the inequality
			\[
				\sum_{i = 1}^{n} \lvert \theta([\mathbf{a}_{i}, \mathbf{b}_{i}], \mathbf{x}) \rvert < 2 \pi
			\]
			implies that $w(\gamma, \mathbf{x}) = 0$.
		\item Let $\mathbf{a}, \mathbf{b}$ be points that lie in a circular disc of radius $r$ and let $\mathbf{x}$ be a point at distance $R > r$ from the centre of the disc. Using a simple trigonometric argument, show:
			\[
				\lvert \theta([\mathbf{a}, \mathbf{b}], \mathbf{x}) \rvert \le 2 \arcsin(r / R)
			\]
		\item Let $\gamma = \sum_{i = 1}^{n} [\mathbf{a}_{i}, \mathbf{b}_{i}]$ be a $1$-cycle such that all the points $\mathbf{a}_{i}, \mathbf{b}_{i}$ are contained in a circular disc of radius $r$. Let $\mathbf{x}$ be a point at distance $R > r$ from the centre of the disc. Suppose
			\[
				R > \frac{r}{\sin{(\pi / n)}}.
			\]
			Combine your results from (i) and (ii) to show that $w(\gamma, \mathbf{x}) = 0$.
		\item Explain why a stronger result than (iii) is totally obvious from the ray-escape formula.
	\end{enumerate}
\end{problem}

\begin{solution}
	\begin{enumerate}
		\item Suppose the inequality holds. Then, by the triangle inequality, the winding number of $\gamma$ about $\mathbf{x}$ satisfies the inequality
			\[
				\lvert w(\gamma, \mathbf{x}) \rvert \le \left \lvert \frac{1}{2 \pi} \sum_{i = 1}^{n} \theta([\mathbf{a}_{i}, \mathbf{b}_{i}], \mathbf{x}) \right \rvert \le \frac{1}{2 \pi} \sum_{i = 1}^{n} \lvert \theta([\mathbf{a}_{i}, \mathbf{b}_{i}], \mathbf{x}) \rvert < 1
			\]
			Since the winding number is an integer, it must be the only integer with norm less than $1$ --- $w(\gamma, \mathbf{x}) = 0$.
		\item Consider the tangents from $\mathbf{x}$ to the circle. These are perpendicular to the radius of the circle at the point. Call the line segment between $\mathbf{x}$ and the centre of the circle $[0, \mathbf{x}]$ and the tangents $[\mathbf{x}, \mathbf{y}_{1}]$ and $[\mathbf{x}, \mathbf{y}_{2}]$.

			Each $[\mathbf{x}, \mathbf{y}_{i}]$ forms a right angled triangle with $[0, \mathbf{x}]$ and $[0, \mathbf{y}_{i}]$. Thus, $\phi = \lvert \theta([0, \mathbf{y}_{i}], \mathbf{x}) \rvert = \arcsin{(r / R)}$. We can see clearly that $\lvert \theta([0, \mathbf{a}], \mathbf{x}) \rvert < \phi$ and similarly for $\mathbf{b}$. Since the circle is convex, we can just add the angles to get the desired inequality ---
			\[
				\lvert \theta([\mathbf{a}, \mathbf{b}], \mathbf{x}) \rvert = \lvert \theta([\mathbf{a}, 0], \mathbf{x}) \rvert + \lvert \theta([0, \mathbf{b}], \mathbf{x}) \rvert < 2 \phi = 2 \arcsin{(r / R)}
			\]
			as desired. Here I haven't been too careful about orientation and sign. I just dded absolute angles since $\mathbf{a}, \mathbf{b}$ lie inside a convex region and we shouldn't run into too many problems. See the picture below.
			\begin{figure}[H]
				\centering
				\includegraphics[width = 0.5 \textwidth]{figures/P2ii circle.pdf}
				\caption{The tangents from $\mathbf{x}$ to the circle subtend an angle $2 \phi$ than $\theta([\mathbf{a}, \mathbf{b}], \mathbf{x})$}
			\end{figure}
	Note that the angle subtended by $[\mathbf{b}, \mathbf{x}]$ and $[0, \mathbf{x}]$ is not generally $\theta([\mathbf{a}, \mathbf{b}], \mathbf{x}) / 2$, but summing with the corresponding angle for $\mathbf{a}$, we do get $\theta([\mathbf{a}, \mathbf{b}], 0)$.
		\item Suppose $R > r / \sin{(\pi / n)}$ or equivalently, $r / R < \sin{(\pi / n)}$. Then since we're in the region where $\sin$ is an increasing function
			\[
				\arcsin{(r / R)} < \arcsin{(\sin{(\pi / n)})} < \pi / n.
			\]

			Since each directed line segment $[\mathbf{a}_{i},  \mathbf{b}_{i}]$ has support inside the circle, it satisfies
			\[
				\lvert \theta([\mathbf{a}_{i}, \mathbf{b}_{i}]) \rvert < 2 \arcsin{(r / R)} < 2 \pi / n
			\]
			Thus,
			\[
				\sum_{i = 1}^{n} \lvert \theta([\mathbf{a}_{i}, \mathbf{b}_{i}], \mathbf{x}) \rvert < 2 \pi,
			\]
			forcing $w(\gamma, \mathbf{x}) = 0$ by part (i).
		\item By the ray escape formula, for any point $\mathbf{x}$ outside the circle, one can pick a ray $R$ in the direction away from the circle. This will have zero crossings with the $1$-cycle, and thus, the $\gamma$, contained in the circle, has winding number $w(\gamma, \mathbf{x}) = 0$ about $\mathbf{x}$.

			In the lecture notes this is generalised to the convex containment formula --- the circle can be any convex set. I think this could be even stronger --- I think this would work for any simply connected set.
	\end{enumerate}
\end{solution}

\pagebreak

\begin{problem}[3.]
	Here are a few of my favourite $1$-cycles. Let $n$ be a positive integer. Define points $v_{\ell}$ on the unit circle by the following formula:
	\[
		v_{\ell} = e^{2 \pi i \ell / n}.
	\]
	Note that these are periodic in that $v_{\ell + n} = v_{\ell}$ for all $\ell$. Let $k$ be any integer. Define the \emph{regular cycle} $E(n, k)$ as follows:
	\[
		E(n, k) = [v_{1}, v_{1 + k}] + [v_{2}, v_{2 + k}] + \dots + [v_{n}, v_{n + k}].
	\]
	In each of the following cases, draw a sketch of the $1$-cycle and write the winding numbers in each component of the complementary region.
	\begin{enumerate}
		\item $E(5, 1)$ 
		\item $E(5, 2)$ 
		\item $E(5, 3)$
		\item $E(6, 3)$ 
		\item $E(8, 2)$ 
		\item $E(8, 3)$ 
		\item $E(2, 1)$
		\item $E(5, -2)$.
		\item Repeat for one more $E(n, k)$ of your choice.
	\end{enumerate}
\end{problem}

\begin{solution}
	In each of the diagrams below, winding numbers of each region are written in blue.
	\begin{enumerate}
		\item The pentagon
			\begin{figure}[H]
				\centering
				\includegraphics[width = 0.5 \textwidth]{figures/P3i E5-1.pdf}
				\caption{$E(5, 1)$}
			\end{figure}
		\item The five point star
			\begin{figure}[H]
				\centering
				\includegraphics[width = 0.5 \textwidth]{figures/P3ii E5-2.pdf}
				\caption{$E(5, 2)$}
			\end{figure}
		\item The reverse (clockwise) five point star
			\begin{figure}[H]
				\centering
				\includegraphics[width = 0.5 \textwidth]{figures/P3iii E5-3.pdf}
				\caption{$E(5, 3)$}
			\end{figure}
		\item The asterisk
			\begin{figure}[H]
				\centering
				\includegraphics[width = 0.5 \textwidth]{figures/P3iv E6-3.pdf}
				\caption{$E(6, 3)$}
			\end{figure}
		\item The double square
			\begin{figure}[H]
				\centering
				\includegraphics[width = 0.5 \textwidth]{figures/P3v E8-2.pdf}
				\caption{$E(8, 2)$}
			\end{figure}
		\item The eight point star
			\begin{figure}[H]
				\centering
				\includegraphics[width = 0.5 \textwidth]{figures/P3vi E8-3.pdf}
				\caption{$E(8, 3)$}
			\end{figure}
		\item The line
			\begin{figure}[H]
				\centering
				\includegraphics[width = 0.5 \textwidth]{figures/P3vii E2-1.pdf}
				\caption{$E(2, 1)$}
			\end{figure}
		\item Also the reverse five point star
			\begin{figure}[H]
				\centering
				\includegraphics[width = 0.5 \textwidth]{figures/P3iii E5-3.pdf}
				\caption{$E(5, -2)$}
			\end{figure}
		\item ``Hexagons are the bestagons'' --- CGP Grey
			\begin{figure}[H]
				\centering
				\includegraphics[width = 0.5 \textwidth]{figures/P3ix E6-1.pdf}
				\caption{$E(6, 1)$}
			\end{figure}
	\end{enumerate}
\end{solution}

\pagebreak

\begin{problem}[4.]
	Obtain a simple general formula for $w(E(n, k), 0)$, under the assumption that $0 < k < n / 2$.
\end{problem}

\begin{solution}
	\[
		\boxed{
			w(E(n, k), 0) = k.
		}
	\]

	Notice that for $k = n / 2$,  $v_{1}$ and $v_{1 + k}$ are an angle of $\pi$ apart, and thus, the line between them passes through the origin. That is, for $k = n / 2$, the winding number about $0$ is not defined.

	Note that
	\begin{align*}
		\operatorname{arg} v_{\ell + k} - \operatorname{arg} v_{\ell} = \left ( \frac{(\ell + k) 2 \pi}{n} + 2 \pi \mathbb{Z} \right ) - \left ( \frac{\ell 2 \pi}{n} + 2 \pi \mathbb{Z} \right ) = \frac{2 \pi k}{n} + 2 \pi \mathbb{Z}.
	\end{align*}
	For $k < n / 2$, this is less than $\pi$ and thus,
	\[
		\theta([v_{\ell}, v_{\ell + k}], 0) = \frac{2 \pi k}{n} < \frac{2 \pi (n / 2)}{n} = \pi.
	\]

	The winding number is the sum of all of these angles ---
	\[
		w(E(n, k), 0) = \frac{1}{2 \pi} \sum_{\ell = 1}^{n} \theta([v_{\ell}, v_{\ell + k}], 0) = \frac{n}{2 \pi} \frac{2 \pi k}{n} = k.
	\]
\end{solution}	

\pagebreak

\begin{problem}[5.]
	Consider the following question: what is the maximum possible value of $w(\gamma, \mathbf{x})$, when $\gamma$ is a $1$-cycle with $n$ edges and $\mathbf{x} \notin \lvert \gamma \rvert$? Let's call that number $w_{\text{max}}(n)$.
	\begin{enumerate}
		\item Let $n$ be a positive integer. Show that $\lvert w(\gamma, \mathbf{x}) \rvert < n / 2$, for all $1$-cycles $\gamma$ with $n$ edges.
		\item For all $n$, exhibit a $1$-cycle $\gamma$ and a point $\mathbf{x} \in \mathbb{R}^{2} \setminus \lvert \gamma \rvert$ such that $w(\gamma, \mathbf{x}) = \lfloor \frac{n - 1}{2} \rfloor$.
	\end{enumerate}
\end{problem}

\begin{solution}
	\begin{enumerate}
		\item By definition, for any $\mathbf{a}, \mathbf{b}, \mathbf{x} \in \mathbb{R}^{2}$, we have $\lvert \theta([\mathbf{a}, \mathbf{b}], \mathbf{x}) \rvert < \pi$ whenever $\mathbf{x} \notin \lvert [\mathbf{a}, \mathbf{b}] \rvert$. Thus, for any $1$-cycle $\gamma = [\mathbf{a}_{1}, \mathbf{b}_{1}] + \dots + [\mathbf{a}_{n}, \mathbf{b}_{n}]$, and any $\mathbf{x} \notin \lvert \gamma \rvert$ we have
			 \begin{align*}
				 \lvert w(\gamma, \mathbf{x}) \rvert & = \left \lvert \frac{1}{2 \pi} \sum_{i = 1}^{n} \theta([\mathbf{a}_{i}, \mathbf{b}_{i}], \mathbf{x}) \right \rvert \\
								     & \le \frac{1}{2 \pi} \sum_{i = 1}^{n} \lvert \theta([\mathbf{a}_{i}, \mathbf{b}_{i}], \mathbf{x}) \rvert \\
								     & < \frac{n \pi}{2 \pi} \\
								     & = \frac{n}{2}.
			\end{align*}
			Here, the first inequality is the triangle inequality and the second follows from the bound on the angle cocycle $\theta$.
		\item
			\[
				\boxed{
					w \left ( E \left ( n, \left \lfloor \frac{n - 1}{2} \right \rfloor \right ), 0 \right ) = \left \lfloor \frac{n - 1}{2} \right \rfloor
			}
			\]
			assuming my answer to problem 4 is correct.
	\end{enumerate}
\end{solution}

\pagebreak

\begin{problem}[6.]
	Let $\gamma$ be a $1$-cycle in $\mathbb{R}^{2}$. We define the \emph{area} of $\gamma$ as follows
	\[
		\operatorname{area}(\gamma) = \int_{\mathbb{R}^{2}} w(\gamma, \mathbf{x}) \, d\mathbf{x}.
	\]
	Equivalently, partition the open set $\mathbb{R}^{2} \setminus \lvert \gamma \rvert$ into its connected components $U_{j}$, and let $w_{j}$ be the value of $w(\gamma, \mathbf{x})$ for $\mathbf{x} \in U_{j}$. Then
	\[
		\operatorname{area}(\gamma) = \sum_{j} w_{j} \operatorname{area}(U_{j}) 
	\]
	where $\operatorname{area}(U_{j})$ is simply the ordinary area of the region $U_{j}$.
	\begin{enumerate}
		\item Compute $\operatorname{area}(\alpha)$, for the (corrected) $1$-cycle $\alpha$ given in question 1.
		\item Construct a quadrilateral $\kappa = [\mathbf{a}, \mathbf{b}] + [\mathbf{b}, \mathbf{c}] + [\mathbf{c}, \mathbf{d}] + [\mathbf{d}, \mathbf{a}]$ with $\operatorname{area}(\kappa) = 0$, such that the four points $\mathbf{a}, \mathbf{b}, \mathbf{c}, \mathbf{d}$ are distinct and do not all lie on a straight line.
	\end{enumerate}
\end{problem}

\begin{solution}
	\begin{enumerate}
		\item
			\[
				\boxed{
					\operatorname{area}(\alpha) = -21/2
				}
			\]

			We can calculate the areas of each of the connected regions of $\alpha$ as follows ---
			\begin{figure}[H]
				\centering
				\includegraphics[width = 0.71 \textwidth]{figures/P2 cycle.pdf}
				\caption{The cycle $\alpha$ with areas in green and winding numbers in blue for each region}
			\end{figure}
			
			Note that we ignored the area of regions with zero winding number since the total area of $\alpha$ is the sum of these areas weighted by winding number. Calculating the total winding number we get
			\[
				\operatorname{area}(\alpha) = (-1)(23 + 5/2 + 1/2) + 1(5/2 + 3/2 + 4 + 3 + 1/2) + 2(2) = -21/2.
			\]
		\item Consider the quadrilateral $\kappa$ with $\mathbf{a} = (-1, 1)$, $\mathbf{b} = (1, -1)$, $\mathbf{c} = (-1, -1)$ and $\mathbf{d} = (1, 1)$. It divides $\mathbb{R}^{2}$ into three connected components --- one unbounded component with winding number $0$ and two bounded components of equal area with winding numbers $1$ and $-1$.

			\begin{figure}[H]
				\centering
				\includegraphics[width = 0.5 \textwidth]{figures/P6ii quad.pdf}
				\caption{The quadrilateral $\kappa$ with winding numbers in blue and areas in green for each (bounded) region}
			\end{figure}
	\end{enumerate}
\end{solution}

\pagebreak

\begin{problem}[7.]
	Give a 1-line proof that $\operatorname{area}(\gamma_{1} + \gamma_{2}) = \operatorname{area}(\gamma_{1}) + \operatorname{area}(\gamma_{2})$ for any two $1$-cycles $\gamma_{1}, \gamma_{2}$.
\end{problem}

\begin{solution}
	\[
		\operatorname{area}(\gamma_{1} + \gamma_{2}) = \int_{\mathbb{R}^{2}} w(\gamma_{1} + \gamma_{2}, \mathbf{x}) \, d\mathbf{x} = \int_{\mathbb{R}^{2}} w(\gamma_{1}, \mathbf{x}) + w(\gamma_{2}, \mathbf{x}) \, d\mathbf{x} = \operatorname{area}(\gamma_{1}) + \operatorname{area}(\gamma_{2}).
	\]
\end{solution}

\pagebreak

\begin{problem}[8.]
	\begin{enumerate}
		\item Let $\gamma = \sum_{i = 1}^{n} [\mathbf{a}_{i}, \mathbf{b}_{i}]$ be a $1$-chain, and let $\mathbf{y} \in \mathbb{R}^{2}$. Show that the $1$-chain
			\[
				\hat{\gamma} = \sum_{i = 1}^{n} ([\mathbf{a}_{i}, \mathbf{b}_{i}] + [\mathbf{b}_{i}, \mathbf{y}] + [\mathbf{y}, \mathbf{a}_{i}])
			\]
			is always a $1$-cycle.
		\item Suppose $\gamma$ itself is a $1$-cycle. Show that $w(\gamma, \mathbf{x}) = w(\hat{\gamma}, \mathbf{x})$ for all $\mathbf{x} \in \mathbb{R}^{2} \setminus \lvert \hat{\gamma} \rvert$.
		\item Deduce that
			\[
				\operatorname{area}(\gamma) = \frac{1}{2} \sum_{i = 1}^{n}  \lVert \mathbf{y} - \mathbf{a}_{i} \rVert \cdot \lVert \mathbf{y} - \mathbf{b}_{i} \rVert \sin{(\theta([\mathbf{a}_{i}, \mathbf{b}_{i}], \mathbf{y}))}.
			\]
	\end{enumerate}
\end{problem}

\begin{solution}
	\begin{enumerate}
		\item Note that each term $[\mathbf{a}_{i}, \mathbf{b}_{i}] + [\mathbf{b}_{i}, \mathbf{y}] + [\mathbf{y}, \mathbf{a}_{i}]$ has each of $\mathbf{a}_{i}$, $\mathbf{b}_{i}$, and $\mathbf{y}$ as the start of a line segment as many times as it is the end of a line segment. Thus, in total, $\mathbf{y}$ and each $\mathbf{a}_{i}$ and $\mathbf{b}_{i}$ are each the start of a line segment as many times as they are the end of a line segment.
		\item By definition,
			\[
				w(\gamma, \mathbf{x}) = \sum_{i = 1}^{n} \theta([\mathbf{a}_{i}, \mathbf{b}_{i}], \mathbf{x}).
			\]
			and
			\begin{align*}
				w(\hat{\gamma}, \mathbf{x}) & = \sum_{i = 1}^{n} \theta([\mathbf{a}_{i}, \mathbf{b}_{i}], \mathbf{x}) + \theta([\mathbf{b}_{i}, \mathbf{y}], \mathbf{x}) + \theta([\mathbf{y}, \mathbf{a}_{i}], \mathbf{x}) \\
						      & = w(\gamma, \mathbf{x}) + \sum_{i = 1}^{n} \theta([\mathbf{b}_{i}, \mathbf{y}], \mathbf{x}) + \theta([\mathbf{y}, \mathbf{a}_{i}], \mathbf{x})
			\end{align*}

			Notice that since $\gamma$ is a $1$-cycle, the sum on the right is the winding number of the sum of $1$-cycles
			\[
				 [\mathbf{a}_{i}, \mathbf{y}] + [\mathbf{y}, \mathbf{b}_{i}] + [\mathbf{b}_{i}, \mathbf{y}] + [\mathbf{y}, \mathbf{a}_{i}].
			\]
			Since $\theta([\mathbf{a}, \mathbf{b}], \mathbf{x}) = - \theta([\mathbf{a}, \mathbf{b}], \mathbf{x})$, each of the four edge $1$-cycles above has winding number $0$.  Thus,
			\[
				\sum_{i = 1}^{n} \theta([\mathbf{b}_{i}, \mathbf{y}], \mathbf{x}) + \theta([\mathbf{y}, \mathbf{a}_{i}], \mathbf{x}) = 0.
			\]
			Substituting this into our expression for $w(\hat{\gamma}, \mathbf{x})$ clearly gives $w(\hat{\gamma}, \mathbf{x}) = w(\gamma, \mathbf{x})$, as desired.
		\item First we calculate the area of $\hat{\gamma}$. 

			We can write $\hat{\gamma}$ as the sum of all triangles $\Delta_{i} = [\mathbf{a}_{i}, \mathbf{b}_{i}] + [\mathbf{b}_{i}, \mathbf{y}] + [\mathbf{y}, \mathbf{a}_{i}]$. A triangle's area is the half the product of two sides and the sine of the angle subtended between them. Accordingly,
			\[
				\operatorname{area}(\Delta_{i}) = \lVert \mathbf{y} - \mathbf{a}_{i} \rVert \lVert \mathbf{y} - \mathbf{b}_{i} \rVert \sin{(\theta, [\mathbf{a}_{i}, \mathbf{b}_{i}], \mathbf{y})}.
			\]
			Since area is additive, the area of $\hat{\gamma}$ is the sum of the areas of all of these triangles. That is,
			\[
				\operatorname{area}(\hat{\gamma}) = \sum_{i = 1}^{n} \operatorname{area}(\Delta_{i}) = \frac{1}{2} \sum_{i = 1}^{n} \lVert \mathbf{y} - \mathbf{a}_{i} \rVert \cdot \lVert \mathbf{y} - \mathbf{b}_{i} \rVert \sin{(\theta([\mathbf{a}_{i}, \mathbf{b}_{i}], \mathbf{y}))}.
			\]

			Since $\hat{\gamma}$ has the same winding number as $\gamma$ everywhere, it must have the same total area (equal functions have equal integrals). Hence, $\operatorname{area}(\gamma) = \operatorname{area}(\hat{\gamma})$ and the desired result follows.
	\end{enumerate}
\end{solution}

\pagebreak

\begin{problem}[9.]
	Calculate $\operatorname{area}(E(n, k))$.
\end{problem}

\begin{solution}
	\[
		\boxed{
			\operatorname{area}(E(n, k)) = \begin{cases}
				n \sin{\left ( \left \lvert \frac{2 \pi k}{n} \right \rvert \right )} & k < n/2 \\
				- n \sin{\left ( \left \lvert \frac{2 \pi k}{n} \right \rvert \right )} & k > n/2 \\
				0
			\end{cases} = \begin{cases}
				n \sin{\left ( \frac{2 \pi k}{n} \right )} & k \neq n / 2 \\
				0 & k = n/2
			\end{cases}
		}
	\]

	For $\gamma = E(n, k)$ and $\mathbf{y} = 0$, consider $\hat{\gamma}$. Now $\hat{\gamma}$ can be seen as the sum of $n$ triangles $\Delta_{i}$. Each of $\Delta_{i}$ has two sides of length $1$ with angle $\lvert 2 \pi k / n \rvert$ between them. Thus, each triangle has $\operatorname{area}(\Delta_{i}) = \sin{(\lvert 2 \pi k / n \rvert)}$. 

	For $k < n / 2$ these triangles are all oriented positively and for $k > n / 2$ these triangles are all oriented negatively. For $k = n / 2$ these triangles are all zero area. Adding up the areas of all of these triangles gives the desired formula.
\end{solution}

\end{document}
